\documentclass[11pt]{article}
\usepackage[margin=1in]{geometry}
\usepackage{hyperref}
\usepackage{enumitem}
\usepackage{booktabs}
\usepackage{xcolor}
\title{ProbOS --- Month 3, Week 14\\Injection Molding Process Qualification (Research Gate First)}
\author{Nisong Monyimba}
\date{\today}

\begin{document}
\maketitle

\section*{Week 14 goal}
Add injection molding process qualification as a new domain model,
reusing the existing, validated engine (\texttt{Distribution}/
\texttt{Model} ABCs, \texttt{MonteCarloEngine}, \texttt{SobolSensitivity},
the now-fixed reproducible seeding, C++ \texttt{LogNormal}/
\texttt{Uniform} classes, PDSL v0.2) --- explicitly NOT a new
architecture, a new model class on infrastructure already at
408/408 green tests.

\textbf{Decision, stated up front:} build the real model first, with
a genuine research gate (Day 1, before any model code), THEN add it
as a fourth tab to the live demo --- not assuming the originating
pitch document's claim that ``public literature on PP/PC process
windows'' exists without verifying it directly.

\section*{Day 1 -- Research gate (before any model code)}
\begin{itemize}[leftmargin=*]
  \item Search for a real, citable, published injection molding
    process-window dataset (PP or PC material) with real reported
    process parameters and real reported out-of-spec/defect outcomes
    --- the exact discipline already used for
    \texttt{PharmaStabilityModel} (Gonzalez-Gonzalez et al. 2023)
  \item \textbf{Explicit decision gate:} if real, usable data is
    found, proceed to Day 2 as planned below. If not found, document
    the gap honestly and either narrow the model's scope to what
    real data DOES support, or defer the week's remaining time to
    another already-planned item
\end{itemize}

\section*{Day 2 -- InjectionMoldingProcessModel (contingent on Day 1)}
\begin{itemize}[leftmargin=*]
  \item Same architectural pattern as \texttt{BatteryModel2Cell}:
    state variables (cavity pressure, melt temperature, fill time,
    cooling time, part dimensional outcome), parameters with real
    uncertainty informed by Day 1's findings (resin lot viscosity ---
    reuses the existing C++ \texttt{LogNormal} class directly; ambient
    humidity and machine repeatability as \texttt{Normal}; wall-
    thickness tolerance from mold wear)
  \item Output: probability of an out-of-spec part as a function of
    the process window --- directly analogous to the battery model's
    thermal-runaway probability output
\end{itemize}

\section*{Day 3 -- Monte Carlo + Sobol sensitivity}
\begin{itemize}[leftmargin=*]
  \item Wire the real \texttt{MonteCarloEngine} and
    \texttt{SobolSensitivity} (now genuinely reproducible, per this
    session's own seeding-bug fix and its permanent regression test)
    to the new model
  \item The actual sellable insight: ``of N process parameters,
    which 2--3 are actually driving scrap rate''
\end{itemize}

\section*{Day 4 -- Validation against the real Day 1 dataset}
\begin{itemize}[leftmargin=*]
  \item Validate against whatever real dataset Day 1 found, with an
    honestly stated tolerance and methodology strength
\end{itemize}

\section*{Day 5 -- PDSL spec for process windows}
\begin{itemize}[leftmargin=*]
  \item A real PDSL v0.2 example letting a process engineer define a
    process window declaratively
\end{itemize}

\section*{Day 6 -- Fourth demo tab}
\begin{itemize}[leftmargin=*]
  \item Add Injection Molding as a fourth tab to the live demo,
    following the EXACT same discipline already established: real
    pre-computed kernel data, no JS toy-formula reimplementation,
    real interpretation bands, real Sobol breakdown, hover tooltips
\end{itemize}

\section*{Day 7 -- Audit, documentation, retrospective}
\begin{itemize}[leftmargin=*]
  \item \texttt{pre\_week\_audit.sh}, full test suite, CI, retrospective
\end{itemize}

\section*{Business framing (for reference, not this week's engineering scope)}
Open-core (Apache 2.0) for the model and Sobol analysis; a paid
audit-trail tier reframes the earlier MoldMind + ValidTrail ideas as
one feature on existing infrastructure. Outreach is explicitly OUT
of this week's scope.

\section*{Exit criteria}
A genuinely validated (or honestly gap-documented) injection molding
domain model, wired to the real Sobol/Monte Carlo engine, with a
working PDSL example and a fourth live demo tab built on real data.


\section*{What was actually accomplished (retrospective)}

\begin{itemize}[leftmargin=*]
  \item A third real, validated domain (\texttt{InjectionMoldingProcessModel})
    added on existing, unmodified infrastructure -- zero new
    architecture, matching the week's own stated goal.
  \item Day 1's research gate found two real, independent sources:
    Kavade \& Kadam (2012) (real physical experiments, the primary
    validation target) and Moayyedian et al. (2021) (FEA simulation,
    used only as a structural cross-check) -- exactly the
    ``don't assume a pitch document's citation claim without
    verifying it'' discipline stated up front in this week's own
    plan.
  \item The model's own real Sobol sensitivity ranking reproduced
    ALL FIVE parameters in the exact same relative order as the
    source paper's independently-measured real ANOVA ranking, now
    covered by a permanent regression test.
  \item An honest architectural difference was documented, not
    hidden: unlike \texttt{BatteryModel2Cell} or
    \texttt{PharmaStabilityModel}, injection molding is a genuinely
    static process (confirmed idempotent) -- the demo tab
    accordingly skips the time-series chart rather than faking a
    flat placeholder trajectory.
\end{itemize}

\section*{Real bugs found and fixed during this week}
\begin{enumerate}[leftmargin=*]
  \item \textbf{PDSL drift semantics (Day 5):} the first PDSL source
    wrote the drift expression as if it directly SET the new state
    value; PDSL's actual, established semantics are Euler-integration
    style (\texttt{new\_state = state + dt*drift}), matching
    Battery/Pharma. Confirmed via cross-validation against the
    hand-written Python model: the mean was off by exactly the
    nominal weight itself (96.826g), diagnosed via direct computation
    before fixing.
  \item \textbf{Readout grid collapse (found via direct user review):}
    hiding grid child elements with \texttt{display:none} does not
    collapse a CSS grid container's own column tracks, leaving a
    large empty blank space. Fixed by dynamically changing
    \texttt{grid-template-columns} on the container itself.
  \item \textbf{Stale Sobol/PDSL panel titles (found via direct user
    review):} both panels were hardcoded to always describe
    ``battery,'' even while displaying injection molding's own real
    data -- confirmed as a real correctness bug, not just a cosmetic
    issue, and fixed with genuine per-tab dynamic content.
\end{enumerate}

\section*{What went right}
\begin{enumerate}[leftmargin=*]
  \item Direct user review of live screenshots caught two real UI
    bugs that automated checks (DOM-ID matching, balance checks)
    could not have caught, since both were about what actually
    RENDERS, not structural correctness.
  \item A suspected third bug (histogram bin coloring) was correctly
    resolved as NOT a bug after full bin-by-bin computation --
    avoided ``fixing'' correctly-working code based on an
    incomplete percentile-level check.
\end{enumerate}

\section*{Numbers at Week 14 close}
\begin{itemize}[leftmargin=*]
  \item Python tests: 419/419 passing
  \item mypy strict: 0 errors: ruff: 0 warnings
  \item Real Sobol structural validation: exact match to the source
    paper's ANOVA ranking, all 5 parameters
  \item Fourth live demo tab, real data, zero known workarounds
\end{itemize}
\end{document}
